\documentclass[10pt,a4paper]{amsart}
\usepackage{amsmath,amssymb,amsthm,graphicx}
% Retained preamble for the Stacks Project chapter {CH} (shared project preamble).
% The chapter file pinned-original-source.tex only carries \input{{preamble}} on its first line,
% so no preamble text lives inside the chapter itself.  The Stacks Project uses one single shared
% preamble for every chapter: the snapshot stages it three times (sources/stacks-chow/preamble.tex,
% sources/stacks-etale-cohomology/preamble.tex and focused-seeds/AGP/sources/stacks-snapshot-20261011/
% preamble.tex) and those copies are identical in size.  The block below is copied mechanically,
% line for line, from sources/stacks-chow/preamble.tex lines 14-18 and 154-217: the Xy-pic setup and
% the native theorem declarations together with the mathematical macros.  Nothing was added, removed
% or redefined; the wrapper class and the amsmath/amssymb/amsthm/graphicx base packages are supplied
% by the assembler.
% For commutative diagrams we use Xy-pic
% Chapter: stacks-cohomology.tex -- Cohomology of Sheaves.
\usepackage[all]{xy}

% We use 2cell for 2-commutative diagrams.
\xyoption{2cell}

% Theorem environments.
%
\theoremstyle{plain}
\newtheorem{theorem}[subsection]{Theorem}
\newtheorem{proposition}[subsection]{Proposition}
\newtheorem{lemma}[subsection]{Lemma}

\theoremstyle{definition}
\newtheorem{definition}[subsection]{Definition}
\newtheorem{example}[subsection]{Example}
\newtheorem{exercise}[subsection]{Exercise}
\newtheorem{situation}[subsection]{Situation}

\theoremstyle{remark}
\newtheorem{remark}[subsection]{Remark}
\newtheorem{remarks}[subsection]{Remarks}

\numberwithin{equation}{subsection}

% Macros
%
\def\lim{\mathop{\mathrm{lim}}\nolimits}
\def\colim{\mathop{\mathrm{colim}}\nolimits}
\def\Spec{\mathop{\mathrm{Spec}}}
\def\Hom{\mathop{\mathrm{Hom}}\nolimits}
\def\Ext{\mathop{\mathrm{Ext}}\nolimits}
\def\SheafHom{\mathop{\mathcal{H}\!\mathit{om}}\nolimits}
\def\SheafExt{\mathop{\mathcal{E}\!\mathit{xt}}\nolimits}
\def\Sch{\mathit{Sch}}
\def\Mor{\mathop{\mathrm{Mor}}\nolimits}
\def\Ob{\mathop{\mathrm{Ob}}\nolimits}
\def\Sh{\mathop{\mathit{Sh}}\nolimits}
\def\NL{\mathop{N\!L}\nolimits}
\def\CH{\mathop{\mathrm{CH}}\nolimits}
\def\proetale{{pro\text{-}\acute{e}tale}}
\def\etale{{\acute{e}tale}}
\def\QCoh{\mathit{QCoh}}
\def\Ker{\mathop{\mathrm{Ker}}}
\def\Im{\mathop{\mathrm{Im}}}
\def\Coker{\mathop{\mathrm{Coker}}}
\def\Coim{\mathop{\mathrm{Coim}}}

% Boxtimes
%
\DeclareMathSymbol{\boxtimes}{\mathbin}{AMSa}{"02}

%
% Macros for moduli stacks/spaces
%
\def\QCohstack{\mathcal{QC}\!\mathit{oh}}
\def\Cohstack{\mathcal{C}\!\mathit{oh}}
\def\Spacesstack{\mathcal{S}\!\mathit{paces}}
\def\Quotfunctor{\mathrm{Quot}}
\def\Hilbfunctor{\mathrm{Hilb}}
\def\Curvesstack{\mathcal{C}\!\mathit{urves}}
\def\Polarizedstack{\mathcal{P}\!\mathit{olarized}}
\def\Complexesstack{\mathcal{C}\!\mathit{omplexes}}
% \Pic is the operator that assigns to X its picard group, usage \Pic(X)
% \Picardstack_{X/B} denotes the Picard stack of X over B
% \Picardfunctor_{X/B} denotes the Picard functor of X over B
\def\Pic{\mathop{\mathrm{Pic}}\nolimits}
\def\Picardstack{\mathcal{P}\!\mathit{ic}}
\def\Picardfunctor{\mathrm{Pic}}
\def\Deformationcategory{\mathcal{D}\!\mathit{ef}}

% --- appended 2026-10-12 (Chief): the project own definitions of the reference/slogan/history
% display environments, copied mechanically from the shared Stacks preamble lines 8-12
% (\usepackage{verbatim} plus the three \newenvironment lines, which the chapter preamble keeps
% verbatim from the project preamble). Nothing added, removed or redefined: these environments
% carry no mathematical content and the project itself typesets them as comments. Without these
% three lines a counted statement whose own text uses them fails with Environment slogan undefined.
\usepackage{verbatim}
\newenvironment{reference}{\comment}{\endcomment}
\newenvironment{slogan}{\comment}{\endcomment}
\newenvironment{history}{\comment}{\endcomment}
\usepackage{cleveref}
\crefname{originallemma}{original lemma}{original lemmas}
\Crefname{originallemma}{Original lemma}{Original lemmas}
\makeatletter
\newlabel{lemma-internal-hom-evaluate}{{0A8U}{source}{}{sourcepointer-1}{}}
\newlabel{lemma-internal-hom-evaluate@cref}{{[originallemma][0][]0A8U}{[1][0][]source}{}{}{}}
\makeatother
\crefname{originallemma}{original lemma}{original lemmas}
\Crefname{originallemma}{Original lemma}{Original lemmas}
\makeatletter
\newlabel{lemma-restriction-RHom-to-U}{{08DL}{source}{}{sourcepointer-2}{}}
\newlabel{lemma-restriction-RHom-to-U@cref}{{[originallemma][0][]08DL}{[1][0][]source}{}{}{}}
\makeatother
\crefname{originalsection}{original section}{original sections}
\Crefname{originalsection}{Original section}{Original sections}
\makeatletter
\newlabel{section-flat}{{06Y7}{source}{}{sourcepointer-3}{}}
\newlabel{section-flat@cref}{{[originalsection][0][]06Y7}{[1][0][]source}{}{}{}}
\makeatother
\crefname{originallemma}{original lemma}{original lemmas}
\Crefname{originallemma}{Original lemma}{Original lemmas}
\makeatletter
\newlabel{lemma-RHom-triangulated}{{08I0}{source}{}{sourcepointer-4}{}}
\newlabel{lemma-RHom-triangulated@cref}{{[originallemma][0][]08I0}{[1][0][]source}{}{}{}}
\makeatother
\crefname{originalreference}{original reference}{original references}
\Crefname{originalreference}{Original reference}{Original references}
\makeatletter
\newlabel{derived-lemma-third-isomorphism-triangle}{{014A}{source}{}{sourcepointer-5}{}}
\newlabel{derived-lemma-third-isomorphism-triangle@cref}{{[originalreference][0][]014A}{[1][0][]source}{}{}{}}
\makeatother
\title{The evaluation map identifies R-Hom(L,M) tensor K with R-Hom(R-Hom(K,L), M) for perfect K}
\author{The Stacks Project Authors}
\date{stacks-project:cohomology.tex:local-snapshot-20261011; GNU FDL 1.2 or later}
\begin{document}
\maketitle
\noindent\textit{Editorial note.} All delivered mathematical text is the author's own, copied verbatim from the pinned Stacks Project chapter file cohomology.tex; no statement or proof step was rewritten, reordered or removed.\par\medskip
\noindent\textit{Editorial note.} The counted span cites six cross-chapter labels, each recorded in assembly.external\_references with its canonical tag as a navigation pointer.\par\medskip
\noindent\textit{Editorial note.} The proof cites no bibliography item inside the counted span.\par\medskip
\noindent\textit{Allowed standard prerequisites (editorial).}\begin{itemize}
\item Behaviour of the internal hom under restriction to opens. Stacks Project, Cohomology of Sheaves, Lemma lemma-restriction-RHom-to-U (tag 0A8M). Justifies the reduction to a local computation on X.
\item The evaluation map and its basic compatibility. Stacks Project, Cohomology of Sheaves, Lemma lemma-internal-hom-evaluate (tag 0A8N). Provides the comparison map whose bijectivity is being proved.
\end{itemize}\medskip
\ifdefined\hypertarget\hypertarget{sourcepointer-1}{}\fi
\noindent\textit{Source reference.} Original lemma 0A8U: sources/stacks-chow/tags.txt, source lines 12927--12927. Cross-chapter standard prerequisite cited by the author inside the counted span; it is used as a pointer and its own proof is not part of this item.\par\medskip
\ifdefined\hypertarget\hypertarget{sourcepointer-2}{}\fi
\noindent\textit{Source reference.} Original lemma 08DL: sources/stacks-chow/tags.txt, source lines 10641--10641. Cross-chapter standard prerequisite cited by the author inside the counted span; it is used as a pointer and its own proof is not part of this item.\par\medskip
\ifdefined\hypertarget\hypertarget{sourcepointer-3}{}\fi
\noindent\textit{Source reference.} Original section 06Y7: sources/stacks-chow/tags.txt, source lines 8876--8876. Cross-chapter standard prerequisite cited by the author inside the counted span; it is used as a pointer and its own proof is not part of this item.\par\medskip
\ifdefined\hypertarget\hypertarget{sourcepointer-4}{}\fi
\noindent\textit{Source reference.} Original lemma 08I0: sources/stacks-chow/tags.txt, source lines 10798--10798. Cross-chapter standard prerequisite cited by the author inside the counted span; it is used as a pointer and its own proof is not part of this item.\par\medskip
\ifdefined\hypertarget\hypertarget{sourcepointer-5}{}\fi
\noindent\textit{Source reference.} Original reference 014A: sources/stacks-chow/tags.txt, source lines 1465--1465. Cross-chapter standard prerequisite cited by the author inside the counted span; it is used as a pointer and its own proof is not part of this item.\par\medskip
\begin{lemma}
\label{lemma-internal-hom-evaluate-isom}
Let $(X, \mathcal{O}_X)$ be a ringed space. Let $K, L, M \in D(\mathcal{O}_X)$.
If $K$ is perfect, then the map
$$
R\SheafHom(L, M) \otimes_{\mathcal{O}_X}^\mathbf{L} K
\longrightarrow
R\SheafHom(R\SheafHom(K, L), M)
$$
of Lemma \ref{lemma-internal-hom-evaluate} is an isomorphism.
\end{lemma}

\begin{proof}
Since the map is globally defined and since formation of the right and
left hand side commute with localization
(see Lemma \ref{lemma-restriction-RHom-to-U}), to prove this we may work
locally on $X$. Thus we may assume $K$ is represented by a strictly
perfect complex $\mathcal{E}^\bullet$.

\medskip\noindent
If $K_1 \to K_2 \to K_3$ is a distinguished triangle in $D(\mathcal{O}_X)$,
then we get distinguished triangles
$$
R\SheafHom(L, M) \otimes_{\mathcal{O}_X}^\mathbf{L} K_1 \to
R\SheafHom(L, M) \otimes_{\mathcal{O}_X}^\mathbf{L} K_2 \to
R\SheafHom(L, M) \otimes_{\mathcal{O}_X}^\mathbf{L} K_3
$$
and
$$
R\SheafHom(R\SheafHom(K_1, L), M) \to
R\SheafHom(R\SheafHom(K_2, L), M)
R\SheafHom(R\SheafHom(K_3, L), M)
$$
See Section \ref{section-flat} and
Lemma \ref{lemma-RHom-triangulated}.
The arrow of Lemma \ref{lemma-internal-hom-evaluate} is functorial in $K$
hence we get a morphism between these distinguished triangles.
Thus, if the result holds for $K_1$ and $K_3$, then the result holds for
$K_2$ by Derived Categories, Lemma
\ref{derived-lemma-third-isomorphism-triangle}.

\medskip\noindent
Combining the remarks above with the distinguished triangles
$$
\sigma_{\geq n}\mathcal{E}^\bullet \to \mathcal{E}^\bullet \to
\sigma_{\leq n - 1}\mathcal{E}^\bullet
$$
of stupid trunctions, we reduce to the case where $K$ consists
of a direct summand of a finite free $\mathcal{O}_X$-module placed
in some degree. By an obvious compatibility of the problem with direct sums
(similar to what was said above) and shifts this reduces us to the case
where $K = \mathcal{O}_X^{\oplus n}$ for some integer $n$.
This case is clear.
\end{proof}

\end{document}
